\section{Generational stratification, quark supports and neutrino realization}
\label{sec:generaciones-quarks-neutrinos-rev2}
\index{fermionic generations!HMT--MD stratification}
\index{quark!cellular supports and record}
\index{PMNS!base transport}

The general law acts on multisections rather than on a list of masses. The
appearance of three generations is therefore expressed as the compatibility of
three structures: the stratification of the atlas, the flavour records, and the
three-dimensional mixing space. In the charged leptonic sector, the
multisections
\begin{equation}
 G_0^{\ell},\qquad G_2^{\ell},\qquad G_4^{\ell}
 \label{eq:tres-multisecciones-leptonicas-rev2}
\end{equation}
have ranks \(1,8,16\). The centre \(G_0^\ell\) contains the rank-one
electronic section; \(G_2^\ell\) and \(G_4^\ell\) instead preserve
nontrivial operator spectra. The complete generational organization
is
\begin{center}
\small
\begin{tabularx}{0.96\textwidth}{@{}c c c c X@{}}
\toprule
\textbf{Generation} & \textbf{Lepton} & \textbf{Par quark} &
\textbf{Sabor neutro} & \textbf{Soporte estructural}\\
\midrule
1 & \(e\) & \((u,d)\) & \(\nu_e\) & centre \(G_0^\ell\), first-generation register and first vector of the interaction basis\\
2 & \(\mu\) & \((c,s)\) & \(\nu_\mu\) & multisection \(G_2^\ell\),
second-generation register and second vector of the interaction basis\\
3 & \(\tau\) & \((t,b)\) & \(\nu_\tau\) & multisection \(G_4^\ell\),
third-generation record, and third vector of the interaction basis\\
\bottomrule
\end{tabularx}
\end{center}
Particle--antiparticle conjugation reverses the charge and orientation
coordinates of the representation but preserves this stratification;
therefore, it does not double the number of generations.

Dimension three also has its own incidence carrier. The oriented
dodecaphasic representation contains
\begin{equation}
 V_{12}\cong V_1\oplus V_3\oplus V_{3'}\oplus V_5,
 \qquad \operatorname{rank}P_3=3,
 \label{eq:descomposicion-rango-tres-generaciones-rev2}
\end{equation}
and the dodecaphasic entangler--Witt transports the positive sector according to
\begin{equation}
 P_G\xleftrightarrow{\;W_{GK}\;}P_3
 \xleftrightarrow{\;U_W\;}Q_3,
 \qquad
 \operatorname{rank}P_G=\operatorname{rank}P_3
 =\operatorname{rank}Q_3=3.
 \label{eq:transporte-rango-tres-generaciones-rev2}
\end{equation}
This three-dimensional carrier permits comparison among generational bases. The
routes, signatures, and mass operators subsequently provide the spectral
content transported by CKM and PMNS.

\subsection{Internal position and records of the \(6\) quark flavors}

The position of a quark flavor is the composite data
\begin{equation}
 \mathfrak q_X=
 (\text{charge branch},\text{generation},(\mathbb Z/3\mathbb Z)_{\rm color},
  \text{multisection},\text{route record}).
 \label{eq:posicion-quark-compuesta-rev2}
\end{equation}
The charge branches \(-1/3\) possess the supports
\begin{align}
 D_1&=\{(4,5),(6,5)\},\nonumber\\
 D_2&=\{(4,3),(4,7),(6,3),(6,7)\},\nonumber\\
 D_3&=\{(2,3),(2,5),(2,7),(8,3),(8,5),(8,7)\},\nonumber\\
 D_4&=\{(2,1),(2,9),(4,1),(4,9),(6,1),(6,9),(8,1),(8,9)\},
 \label{eq:soportes-quark-down-rev2}
\end{align}
whereas the charge branches \(+2/3\) occupy
\begin{align}
 U_1&=\{(4,4),(4,6),(6,4),(6,6)\},\nonumber\\
 U_3&=\{(2,2),(2,4),(2,6),(2,8),(4,2),(4,8),
       (6,2),(6,8),(8,2),(8,4),(8,6),(8,8)\}.
 \label{eq:soportes-quark-up-rev2}
\end{align}
Each support lifts to four orientations per cell. The colour quotient
organizes \(36=12_{\rm R}+12_{\rm G}+12_{\rm B}\) oriented biographies and preserves the same
mass per flavour before the colour realization.

The nominal transducer
\begin{equation}
 \mathcal N_{\rm ruta}(\Gamma_f)
 =((N,E,S,O),h_{369},r_5)
 \in\mathbb Z_{\ge0}^{4}\times\mathbb Z_{\ge0}^{2}
 \label{eq:transductor-nominal-quark-rev2}
\end{equation}
publish the following coordinates before any scalar realization:
\begin{center}
\small
\begin{tabularx}{0.96\textwidth}{@{}c c c c c X@{}}
\toprule
\textbf{Flavour} & \textbf{Generation} & \textbf{Branch} &
\textbf{\((N,E,S,O)\)} & \textbf{\((h_{369},r_5)\)} &
\textbf{Eje dominante}\\
\midrule
\(u\)&1&\(+2/3\)&\((34,31,31,12)\)&\((51,13)\)&N--S\\
\(d\)&1&\(-1/3\)&\((27,35,18,28)\)&\((47,17)\)&E--O\\
\(c\)&2&\(+2/3\)&\((20,38,24,26)\)&\((49,16)\)&E--O\\
\(s\)&2&\(-1/3\)&\((50,32,12,14)\)&\((57,11)\)&N--S\\
\(t\)&3&\(+2/3\)&\((46,27,16,19)\)&\((57,9)\)&N--S\\
\(b\)&3&\(-1/3\)&\((12,63,9,24)\)&\((47,7)\)&E--O\\
\bottomrule
\end{tabularx}
\end{center}
The associated nominal signatures are
\begin{center}
\small
\begin{tabular}{@{}c c c@{}}
\toprule
\textbf{Flavor} & \textbf{\((n_A,s_C)\)} & \textbf{\((W_+,W_-)\)}\\
\midrule
\(u\)&\((11,7)\)&\((73,59)\)\\
\(d\)&\((17,8)\)&\((110,94)\)\\
\(c\)&\((61,8)\)&\((374,358)\)\\
\(s\)&\((41,-3)\)&\((243,249)\)\\
\(t\)&\((100,-1)\)&\((599,601)\)\\
\(b\)&\((71,-5)\)&\((421,431)\)\\
\bottomrule
\end{tabular}
\end{center}
These coordinates distinguen flavors that share branch or generation and close
the morphism of flavor without consulting no mass. For a section physical (P),
we define
\begin{equation}
 \operatorname{sabor}(P)
 =\bigl(\sigma_{\rm ud}(P),g(P),\chi_{\rm fina}(P)\bigr),
 \qquad
 \chi_{\rm fina}(P)
 =(\text{rotations},\text{switches},\text{axis},N,E,S,O,h_{369},r_5),
 \label{eq:morfismo-sabor-ruta-cerrado-rev2}
\end{equation}
and, for quarks, the ternary orbit
\([\gamma_P]_{\rm col}=\{\gamma_P^R,\gamma_P^G,\gamma_P^B\}\) is also preserved. The physical
realization is the map already determined by the preceding data,
\begin{equation}
 \mathcal S_{\rm phys}(P,\mathfrak M):
 \bigl(\sigma_{\rm ud},g,\chi_{\rm fina},[\gamma_P]_{\rm col}\bigr)
 \longmapsto \gamma^{\rm surv}_{P,\mathfrak M}
 \longmapsto
 \bigl(\sigma(\gamma^{\rm surv}_{P,\mathfrak M}),
       B_F^D,B_F^\Omega\bigr),
 \label{eq:realizacion-fisica-sabor-ruta-cerrada-rev2}
\end{equation}
where \(\mathfrak M\) is the compatible coupling and
\(\gamma^{\rm surv}_{P,\mathfrak M}\) the route that persists under it. For
charged leptons the color orbit is reduced to the terminal object. The mass
character is then evaluated on that same genealogy:
\begin{equation}
 \boxed{
 \frac{m_{\mathfrak M}(P)}{m_e^{\rm HMT}}
 =\chi_{\rm mass}\!\left(
   \sigma\!\left(\gamma^{\rm surv}_{P,\mathfrak M}\right)\right).}
 \label{eq:masa-como-caracter-persistencia-rev2}
\end{equation}
Therefore, flavor, generation, and color do not await a future selector: they publish the
route class and its eigenvector before opening the external comparison. Mass does not
retrospectively choose the route; it only reads multiplicatively the persistence
that the coupling has stabilized.

\subsection{Fermionicity, mass eigenvalues, and PMNS transport}

The exterior closure of the exchange functor constructs the CAR relations
and fixes fermionicity. In the neutral sector, the label
\(\nu_e,\nu_\mu,\nu_\tau\) designates an interaction basis; the mass
values are eigenvalues of a mixed neutral operator. The internal frames
\(F_\ell,F_\nu:\mathbb C^3\to\operatorname{im}P_3\) compare those bases, and
the matrix
\begin{equation}
 U_{\rm HMT}=F_\ell^*\Phi_LF_\nu
 \label{eq:transporte-pmns-generaciones-rev2}
\end{equation}
is a unitary transport. PMNS changes basis between interaction states
and mass eigenvectors; the neutral scale belongs to the operator spectrum,
not to the change of basis.

\subsection{Neutral extractor Z0 prior to projection}
\label{sec:extractor-neutro-z0-torsion-rev2}

The recovered construction of the neutral sector acts on the enriched
dodecaphase record, before collapsing sheet, torsion, and the
bidirectional bit. Let \(E_m\), \(m\in\mathbb Z/12\mathbb Z\) be the twelve
nine-tick events of the CT--108 record. At each tick \(t\), the following are preserved
\[
 \theta(t)=\mathbf1_{\{d(t)=9\}},\qquad
 \varepsilon(t)\in\{-1,+1\},\qquad
 \kappa(t)=\mathbf1_{\{9\ {\rm en\ corona}\}},
 \qquad \omega(t)=(-1)^{\kappa(t)}.
\]
The first factor detecta the torsion nonadic, the second preserves the sheet
two--ways and the third changes the sign in the turn of \(180^\circ\) of the
crown. The integer torsional of each event and its difference antipodal are
\begin{equation}
 T_\nu(m)=\sum_{t\in E_m}
   \omega(t)\varepsilon(t)\theta(t),\qquad
 e_m^\nu=T_\nu(m)-T_\nu(m+6),
 \label{eq:extractor-z0-torsion-antipodal-rev2}
\end{equation}
and the neutral trit is fixed, without continuous parameters, by
\begin{equation}
 \tau_m^\nu=\operatorname{sgn}(e_m^\nu)
 \in\{-1,0,+1\}.
 \label{eq:trit-neutro-torsion-antipodal-rev2}
\end{equation}

For the four jump orbits four
\[
 I_a=\{a,a+4,a+8\}\pmod{12},\qquad a=0,1,2,3,
\]
the event coefficient entering into the signature and the majority sign observable are, respectively,
\begin{equation}
 b_{120}^\nu
 =\sum_{a=0}^{3}\operatorname{sgn}
   \left(\sum_{m\in I_a}e_m^\nu\right),
 \qquad
 \nu_{120}^{\nu,\mathrm{sgn}}
 =\sum_{a=0}^{3}\operatorname{sgn}
   \left(\sum_{m\in I_a}\tau_m^\nu\right).
 \label{eq:canales-120-neutros-distintos-rev2}
\end{equation}
These two integers are not identified: the first is the coordinate
\(b_{120}=\nu_{120}^{\rm ev}\) of the character, while the second is a
diagnostic of majority. The quadratic memory is
\begin{equation}
 b_\zeta^\nu=\nu_{270}^\nu
 =\sum_{m=0}^{11}\tau_m^\nu\tau_{m+3}^\nu.
 \label{eq:canal-270-neutro-rev2}
\end{equation}
Thus, for each surviving neutral history
\(\gamma_{\nu_i}\), the register publishes, before any comparison, the signature
\[
 v_\nu(\gamma_{\nu_i})
 =\bigl(0,0,k_{\Delta_4}^{(i)},b_{120}^{\nu_i},
        b_\zeta^{\nu_i}\bigr)
\]
and the relative exponent
\begin{equation}
 \mathcal A_\nu^{(i)}
 =k_{\Delta_4}^{(i)}\Delta_4
  +b_{120}^{\nu_i}s_{120}
  +b_\zeta^{\nu_i}\zeta_{270},
 \qquad
 \zeta_{270}=135\Delta_4^2.
 \label{eq:caracter-relativo-neutro-z0-rev2}
\end{equation}
Here \(\zeta_{270}\) is not replaced by the tower moment
\(s_{270}=q_-^{270}-q_+^{270}\).

The global two-way conjugation
\(\tau^\nu\mapsto-\tau^\nu\) reverses
\(b_{120}^\nu\) and preserves \(b_\zeta^\nu\). Consequently, it generates two
branches of the relative spectrum,
\begin{equation}
 \mathcal A_{\nu,\pm}^{(i)}
 =k_{\Delta_4}^{(i)}\Delta_4
  \pm b_{120}^{\nu_i}s_{120}
  +b_\zeta^{\nu_i}\zeta_{270},
 \qquad
 \text{NO}\longleftrightarrow\text{IO},
 \label{eq:bifurcacion-neutra-no-io-rev2}
\end{equation}
which are distinguished by the subsequent sign of \(\Delta m_{31}^2\), not by a selection using PDG data. This construction closes the discrete Z0 extractor and the relative bifurcation. The common absolute scale \(m_0\) and the full-rank physical PMNS kernel remain separate objects: they are neither introduced into the selector nor declared to be deduced from two--ways parity alone.

The realization with a single phase of rank one produces
\[
 (\theta_{12},\theta_{23},\theta_{13})
 =(19.6807^\circ,56.4970^\circ,29.6224^\circ)
\]
and does not reproduce mixing phenomenology. This control selects, as the appropriate domain, a
full-rank register kernel,
\begin{equation}
 G_{ab}=\frac1{\sqrt{|L_a||N_b|}}
 \sum_{c\in L_a}\sum_{d\in N_b}
 K_{\rm registro}(c,d;\xi_c,\xi_d),
 \qquad U_{\rm int}=\operatorname{polar}(G),
 \label{eq:kernel-registro-rango-completo-rev2}
\end{equation}
whose technical source is developed in the chapter on CKM and PMNS
transport. The mass operator, its eigenvalues, and the flavor-reading map
are thereby kept separate.
